\documentclass{article}
\usepackage{amsmath,amssymb,xurl}
\usepackage[hidelinks]{hyperref}
% Shared commands for notes in docs/paper-gaps/.

\DeclareUnicodeCharacter{2080}{\ifmmode{}_{0}\else\textsubscript{0}\fi}
\DeclareUnicodeCharacter{2081}{\ifmmode{}_{1}\else\textsubscript{1}\fi}
\DeclareUnicodeCharacter{2082}{\ifmmode{}_{2}\else\textsubscript{2}\fi}
\DeclareUnicodeCharacter{2099}{\ifmmode{}_{n}\else\textsubscript{n}\fi}

\newcommand{\Lean}{\textsc{Lean}}
\ExplSyntaxOn
\NewDocumentCommand{\leanid}{m}{
  \ifmmode
    \textsf{\detokenize{#1}}
  \else
    \group_begin:
    \urlstyle{sf}
    \tl_set:Nn \l_tmpa_tl {#1}
    \tl_replace_all:Nnn \l_tmpa_tl {\_} {_}
    \exp_args:NV \path \l_tmpa_tl
    \group_end:
  \fi
}
\ExplSyntaxOff
\providecommand{\tr}{\operatorname{tr}}
\newcommand{\tnleanIssueBase}{https://github.com/LionSR/TNLean/issues/}
\newcommand{\tnleanPrBase}{https://github.com/LionSR/TNLean/pull/}
\newcommand{\tnleanDocBase}{https://sirui-lu.com/TNLean/docs/}
\newcommand{\ghissue}[1]{\href{\tnleanIssueBase#1}{GitHub issue \##1}}
\newcommand{\ghpr}[1]{\href{\tnleanPrBase#1}{PR \##1}}
\newcommand{\doclink}[1]{\href{\tnleanDocBase#1.html}{\path{#1}}}

% Dirac notation and the identity operator, as in blueprint/src/macros/common.tex.
% \providecommand keeps note-local definitions authoritative.
\usepackage{dsfont}
\providecommand{\ket}[1]{|#1\rangle}
\providecommand{\bra}[1]{\langle #1|}
\providecommand{\braket}[2]{\langle #1 | #2 \rangle}
\providecommand{\ketbra}[2]{|#1\rangle\!\langle #2|}
\providecommand{\Id}{\mathds{1}}
\providecommand{\one}{\mathds{1}}
\providecommand{\C}{\mathbb{C}}
\providecommand{\MN}[1]{M_{#1}(\C)}
\providecommand{\MD}{M_D(\C)}

% Verdict marker: \gapnote{<kind>}{<status>}. Kind is the policy
% classification (clarification, local-correction, scope-restriction,
% unfaithful, false-source, open-gap); status is open, wip (elimination
% underway), resolved, or historical. Machine-read by the site index;
% prints a verdict line.
\newcommand{\gapnote}[2]{\par\noindent\textbf{Verdict:} #1 (#2).\par}

\title{An Explicit Finite Presentation for Restricted Native Dual Homotopy}
\author{}
\date{}
\begin{document}
\maketitle
\gapnote{scope-restriction}{open}

For the generated supported-homotopy reading of Schuch, Cirac, and
P\'erez-Garc\'ia, \href{https://arxiv.org/abs/1001.3807v3}{arXiv:1001.3807v3},
Lemma 6.14, source \path{Papers/1001.3807/paper_v3.tex}, lines 2199--2214,
one can replace quantification over every flat fibre transport by equality
in one explicit finitely presented group.

The native positive east and north steps give two generators at each
actual torus vertex. The reverse steps use inverse letters and path words
are reverse-ordered. The relators are exactly the actual square words at
the allowed primal supports and the actual words of fixed native comb
routes from the origin to each vertex. There are at most twice as many
relators as vertices. Neither the generators nor these local geometric
relations refer to the path pair to be compared or to its homotopy class.
The finite relation set is
\leanid{TNLean.PEPS.torusDualPresentationRelators}.

The theorem
\leanid{TNLean.PEPS.torusDualHomotopy_iff_presentation_word_eq}
proves that two same-endpoint native paths are supported-homotopic exactly
when their actual words have equal image in this presentation. The converse
normalizes an arbitrary locally flat transport by the fixed root routes.
Its normalized step permutations satisfy every relator, so the presentation
acts on the root fibre. Equal group words give equal original transport
actions, and the earlier semantic completeness theorem gives homotopy.
No homotopy witness or equality of the desired path classes is supplied.

Both periods are positive, making the native vertex and generator sets
finite. Periods one and two are included with all native direction labels
retained. The earlier semantic classification at period zero is unchanged.
This is a finite-presentation classification, not a terminating word
algorithm: equality in the presented group remains the algebraic criterion.
A finite geometric normal form, complexity bounds, and a finite-group test
family are not proved here. Displacement alone still loses the nonabelian
obstruction whenever a primal square support is missing.

This homotopy result is distinct from physical vacuum-density admissibility
and from the annular physical Wilson-loop obstruction. The broader geometry
and algorithmic classification remain open under \ghissue{8703}; no
\ghissue{8947} Hamiltonian files or claims are changed.
\end{document}
